\scanpage{5}
% archive foliation: 4
Donc: se donner un morphisme de topos \inserted{$\mathcal X_0\mathbin{-}\mathcal X_1$} revient au $=$ que de se donner un topos $\mathcal X$, \inserted{\unsure{distinguer un}} un sous-topos ouvert \inserted{$\mathcal X_0$}, et le fermé complémentaire $\mathcal X_1$, de telle façon que le foncteur recollement $\varphi$ soit, non seulement exact à g, mais commute aux $\underset{\leftarrow}{\lim}$ quelconques, \deleted{\unsure{et de}} i.e. admette un adjoint à gauche $f^*$, et que de \uline{plus} ce dernier soit exact à gauche. Alors l'inclusion $i:\mathcal X_1\longrightarrow\mathcal X$ est fort asphérique (\uline{réciproque} ?) et la suite exacte de cohomologie relative
\margin{?}
\[
(5)\quad\longrightarrow H^i_{\mathcal X_1}(\mathcal X,F)\longrightarrow H^i(\mathcal X,F)\longrightarrow H^i(\mathcal X_0,F)\longrightarrow H^{i+1}_{\mathcal X_{\unsure{1|0}}}(\mathcal X,F)\longrightarrow\cdots
\]
peut s'écrire
\[
(5')\quad\longrightarrow H^i_{\mathcal X_1}(\mathcal X,F)\longrightarrow H^i(\mathcal X_1,F_1)\longrightarrow H^i(\mathcal X_0,F_0)\longrightarrow H^{i+1}_{\mathcal X_{\unsure{1|0}}}(\mathcal X,F)\longrightarrow\cdots
\]
On \unsure{pose}
\[
H^i_{\mathcal X_1}(\mathcal X,F)\overset{\text{déf}}{=}H^i(\underbrace{\mathcal X_1\operatorname{mod}\mathcal X_0}_{\inserted{f\ \text{sous-entendu}}},F=(F_0,F_1,u))
\]
noté aussi \deleted{\unsure{en}} $H^i(\mathcal X_1\operatorname{mod}\mathcal X_0\unsure{\}},F_1)$ si $F_0=f^*(F_1)$, $u$ tautolog.
